% GENERATED FILE. Edit canonical lessons, then run npm run generate:books.
% canonical-source-hash: fnv1a64:c08c0fd1f570f425
% canonical-lessons: FA-C426-rahbandan, FA-C426-park-kardan, FA-C426-dor-zadan, FA-C426-yektarafe, FA-C426-govahiname

\chapter{Traffic jams and parking}
\label{ch:fa-traffic-jams-and-parking}
\hlchaptermodality{426}

\begin{chapteropening}
\textbf{By the end of this chapter:} \emph{I can talk about jams and parking: a traffic jam, parking, turning round, one-way streets and the driving licence.}

Complete the last lesson of chapter 426: I can talk about jams and parking: a traffic jam, parking, turning round, one-way streets and the driving licence.
\end{chapteropening}

\section[\texorpdfstring{\fa{راهبندان} (râhbandân)}{راهبندان (râhbandân)}]{\fa{راهبندان} (râhbandân) — a traffic jam, a hold-up}

\label{lesson:FA-C426-rahbandan}



\begin{quote}
Before the new one: say the Persian for traffic, then the Persian for to swerve.
\end{quote}

\subsection*{You'll want to know: \fa{راهبندان}}
\textbf{\fa{راهبندان}} (\emph{râhbandân}) — \textquotedblleft{}a traffic jam, a hold-up\textquotedblright{}.

\begin{quote}
\fa{به} \fa{خاطر} \fa{راهبندان} \fa{دیر} \fa{به} \fa{خانه} \fa{رسیدیم}.
\end{quote}

\begin{quote}
\emph{Be khâter-e râhbandân dir be khâne residim.}
\end{quote}

\begin{quote}
Because of the traffic jam we got home late.
\end{quote}

Literally 'road-blocking': \emph{râh} 'road' + \emph{band} 'tie, block' + \emph{-ân}.

\begin{cousinweb}
\emph{râhbandân} is a native compound: \emph{râh} 'road' goes back to Middle Persian \emph{rāh} (EtymDB), and \emph{band} 'a tie, a block' to the root of \emph{bastan} 'to bind', which EtymDB traces to PIE \emph{b\textsuperscript{h}end\textsuperscript{h}-}.
\end{cousinweb}

\begin{grammarlens}[title={What you've built}]
One more word for telling, explaining and describing at B1.
\end{grammarlens}

\subsection*{Your turn}
\begin{itemize}
\raggedright
  \item \emph{Say it:} \emph{râhbandân}
  \item \emph{Say it:} \emph{râhbandân}, once more
  \item \emph{Say it:} a sentence of your own with \emph{râhbandân}
  \item \emph{Recall:} say the Persian for traffic, then the Persian for to swerve, then say \emph{râhbandân} again
\end{itemize}

\begin{tcolorbox}[breakable,colback=teal!4,colframe=teal!35!black,title={Before you move on}]
What does \fa{راهبندان} mean? (\textquotedblleft{}A traffic jam, a hold-up\textquotedblright{}.) Say it once more.
\end{tcolorbox}
\section[\texorpdfstring{\fa{پارک} \fa{کردن} (pârk kardan)}{پارک کردن (pârk kardan)}]{\fa{پارک} \fa{کردن} (pârk kardan) — to park}

\label{lesson:FA-C426-park-kardan}



\begin{quote}
Before the new one: say the Persian for to swerve, then the Persian for a traffic jam.
\end{quote}

\subsection*{You'll want to know: \fa{پارک} \fa{کردن}}
\textbf{\fa{پارک} \fa{کردن}} (\emph{pârk kardan}) — \textquotedblleft{}to park\textquotedblright{}.

\begin{quote}
\fa{جلوی} \fa{بیمارستان} \fa{نمیشود} \fa{پارک} \fa{کرد}.
\end{quote}

\begin{quote}
\emph{Jolow-ye bimârestân nemishavad pârk kard.}
\end{quote}

\begin{quote}
You cannot park in front of the hospital.
\end{quote}

\emph{pârk} + \emph{kardan}. Present stem \emph{kon-}: \emph{pârk mikonam}.

\begin{cousinweb}
\emph{pârk kardan} is built from \emph{pârk}, a loan usually traced to French \emph{parc} 'park', and \emph{kardan} 'to do', inherited from Middle Persian \emph{kardan} (EtymDB).
\end{cousinweb}

\begin{grammarlens}[title={What you've built}]
One more word for telling, explaining and describing at B1.
\end{grammarlens}

\subsection*{Your turn}
\begin{itemize}
\raggedright
  \item \emph{Say it:} \emph{pârk kardan}
  \item \emph{Say it:} \emph{pârk kardan}, once more
  \item \emph{Say it:} a sentence of your own with \emph{pârk kardan}
  \item \emph{Recall:} say the Persian for to swerve, then the Persian for a traffic jam, then say \emph{pârk kardan} again
\end{itemize}

\begin{tcolorbox}[breakable,colback=teal!4,colframe=teal!35!black,title={Before you move on}]
What does \fa{پارک} \fa{کردن} mean? (\textquotedblleft{}To park\textquotedblright{}.) Say it once more.
\end{tcolorbox}
\section[\texorpdfstring{\fa{دور} \fa{زدن} (dowr zadan)}{دور زدن (dowr zadan)}]{\fa{دور} \fa{زدن} (dowr zadan) — to turn round, to make a U-turn; to go round}

\label{lesson:FA-C426-dor-zadan}



\begin{quote}
Before the new one: say the Persian for a traffic jam, then the Persian for to park.
\end{quote}

\subsection*{You'll want to know: \fa{دور} \fa{زدن}}
\textbf{\fa{دور} \fa{زدن}} (\emph{dowr zadan}) — \textquotedblleft{}to turn round, to make a U-turn; to go round\textquotedblright{}.

\begin{quote}
\fa{راه} \fa{را} \fa{اشتباه} \fa{رفتیم} \fa{و} \fa{دور} \fa{زدیم}.
\end{quote}

\begin{quote}
\emph{Râh râ eshtebâh raftim va dowr zadim.}
\end{quote}

\begin{quote}
We went the wrong way and turned round.
\end{quote}

\emph{dowr} 'turn, round' + \emph{zadan}. \emph{Dowr-e meydân} — round the square.

\begin{cousinweb}
\emph{dowr zadan} is built from \emph{dowr} 'a turn, a round', borrowed from Arabic \emph{dawr} 'turn' (EtymDB), and the native verb \emph{zadan} 'to strike'.
\end{cousinweb}

\begin{grammarlens}[title={What you've built}]
One more word for telling, explaining and describing at B1.
\end{grammarlens}

\subsection*{Your turn}
\begin{itemize}
\raggedright
  \item \emph{Say it:} \emph{dowr zadan}
  \item \emph{Say it:} \emph{dowr zadan}, once more
  \item \emph{Say it:} a sentence of your own with \emph{dowr zadan}
  \item \emph{Recall:} say the Persian for a traffic jam, then the Persian for to park, then say \emph{dowr zadan} again
\end{itemize}

\begin{tcolorbox}[breakable,colback=teal!4,colframe=teal!35!black,title={Before you move on}]
What does \fa{دور} \fa{زدن} mean? (\textquotedblleft{}To turn round, to make a U-turn; to go round\textquotedblright{}.) Say it once more.
\end{tcolorbox}
\section[\texorpdfstring{\fa{یک} \fa{طرفه} (yektarafe)}{یک طرفه (yektarafe)}]{\fa{یک} \fa{طرفه} (yektarafe) — one-way (of a street)}

\label{lesson:FA-C426-yektarafe}



\begin{quote}
Before the new one: say the Persian for to park, then the Persian for to turn round.
\end{quote}

\subsection*{You'll want to know: \fa{یک} \fa{طرفه}}
\textbf{\fa{یک} \fa{طرفه}} (\emph{yektarafe}) — \textquotedblleft{}one-way (of a street)\textquotedblright{}.

\begin{quote}
\fa{این} \fa{کوچه} \fa{یک} \fa{طرفه} \fa{است}.
\end{quote}

\begin{quote}
\emph{In kuche yektarafe ast.}
\end{quote}

\begin{quote}
This lane is one-way.
\end{quote}

\emph{yek} 'one' + \emph{taraf} 'side' + \emph{-e}. A two-way street is \emph{dotarafe}.

\begin{cousinweb}
\emph{yektarafe} is built from native \emph{yek} 'one' and \emph{taraf} 'side', a loan from Arabic \emph{ṭaraf} 'side, edge', with the ending \emph{-e}.
\end{cousinweb}

\begin{grammarlens}[title={What you've built}]
One more word for telling, explaining and describing at B1.
\end{grammarlens}

\subsection*{Your turn}
\begin{itemize}
\raggedright
  \item \emph{Say it:} \emph{yektarafe}
  \item \emph{Say it:} \emph{yektarafe}, once more
  \item \emph{Say it:} a sentence of your own with \emph{yektarafe}
  \item \emph{Recall:} say the Persian for to park, then the Persian for to turn round, then say \emph{yektarafe} again
\end{itemize}

\begin{tcolorbox}[breakable,colback=teal!4,colframe=teal!35!black,title={Before you move on}]
What does \fa{یک} \fa{طرفه} mean? (\textquotedblleft{}One-way (of a street)\textquotedblright{}.) Say it once more.
\end{tcolorbox}
\section[\texorpdfstring{\fa{گواهینامه} (govâhinâme)}{گواهینامه (govâhinâme)}]{\fa{گواهینامه} (govâhinâme) — a driving licence; a certificate}

\label{lesson:FA-C426-govahiname}



\begin{quote}
Before the new one: say the Persian for to turn round, then the Persian for one-way.
\end{quote}

\subsection*{You'll want to know: \fa{گواهینامه}}
\textbf{\fa{گواهینامه}} (\emph{govâhinâme}) — \textquotedblleft{}a driving licence; a certificate\textquotedblright{}.

\begin{quote}
\fa{برادرم} \fa{پارسال} \fa{گواهینامه} \fa{گرفت}.
\end{quote}

\begin{quote}
\emph{Barâdaram pârsâl govâhinâme gereft.}
\end{quote}

\begin{quote}
My brother got his driving licence last year.
\end{quote}

\emph{govâhi} 'witness, proof' + \emph{nâme} 'letter'. In full: \emph{govâhinâme-ye rânandegi}.

\begin{cousinweb}
\emph{govâhinâme} is a native compound of \emph{govâhi} 'witness, proof', from \emph{govâh} 'a witness', and \emph{nâme} 'letter, document', usually traced to Middle Persian \emph{nāmag}.
\end{cousinweb}

\begin{grammarlens}[title={What you've built}]
One more word for telling, explaining and describing at B1.
\end{grammarlens}

\subsection*{Your turn}
\begin{itemize}
\raggedright
  \item \emph{Say it:} \emph{govâhinâme}
  \item \emph{Say it:} \emph{govâhinâme}, once more
  \item \emph{Say it:} a sentence of your own with \emph{govâhinâme}
  \item \emph{Recall:} say the Persian for to turn round, then the Persian for one-way, then say \emph{govâhinâme} again
\end{itemize}

\begin{tcolorbox}[breakable,colback=teal!4,colframe=teal!35!black,title={Before you move on}]
What does \fa{گواهینامه} mean? (\textquotedblleft{}A driving licence; a certificate\textquotedblright{}.) Say it once more.
\end{tcolorbox}
